\chapter{Когда писать, а когда читать: добавляем \nt{ACET}}
\chaptermark{Добавляем \nt{ACET}}
\label{sec:acet}

\section{Чего не хватает}

У микросхемы памяти, кроме проводов адреса и данных, есть ещё одна линия --- разрешение записи. Пока она выключена, память только читается; когда включена, то, что лежит на линиях данных, записывается. Без такой линии память не работает: если каждое чтение одновременно что-то записывает, прочитанное, в том числе прочитанное с ошибкой, тут же записывается обратно.

В мозге эта задача стоит острее, чем в микросхеме. В главе~\ref{sec:glutcomputer} память была группой \nt{GLUT}-нейронов, которую держат включённой её собственные связи (рис.~\ref{fig:bistable}). В главе~\ref{sec:dofa} мы видели, что эти связи учатся правилом Хебба прямо по ходу работы. Значит, одни и те же синапсы и хранят старое, и принимают новое, а отдельной фазы записи, как и часов, нет. Чем сеть отличает момент, когда нужно запомнить то, что видишь, от момента, когда нужно вспомнить то, что знаешь? В главе~\ref{sec:neurontypes} мы предположили, что эту линию держит \nt{ACET}. В этой главе мы проверим, как такая линия должна работать и почему без неё не обойтись.

\section{Три действия одного провода}

Ацетилхолиновых нейронов, работающих внутри мозга, немного, и собраны они в нескольких небольших группах, а их выходы расходятся по всей коре. Это широковещательный канал, такой же, как \nt{DOPA} и \nt{NORA}. Уровень ацетилхолина в коре высок при внимательном бодрствовании и низок в медленном сне~\cite{Hasselmo1999}. Как и у дофамина, смысл сигнала задаёт рецептор получателя (утверждение~\ref{prop:receptor}): у ацетилхолина есть и быстрые рецепторы, открывающие канал, и медленные, запускающие реакции внутри клетки. Для нас важны три действия.

\emph{Первое: ослабить внутренние связи.} Опыты на срезах обонятельной коры показали, что ацетилхолин сильно подавляет передачу по связям между нейронами одной области и почти не трогает входы, приходящие в неё извне~\cite{HasselmoBower1992}.

\emph{Второе: усилить входы.} Ацетилхолин повышает возбудимость нейронов и облегчает передачу по некоторым входным путям; сигнал от органов чувств становится громче собственной активности сети~\cite{Hasselmo2006}.

\emph{Третье: облегчить изменение весов.} При высоком уровне ацетилхолина связи меняются легче~\cite{Hasselmo2006}.

Для инженера все три действия --- одна операция: переключение из режима «читать» в режим «писать». Сеть перестаёт слушать саму себя, начинает слушать вход и запоминает то, что слышит.

\section{Сеть, которая дописывает}

Возьмём $N$ \nt{GLUT}-нейронов, соединённых друг с другом. Сохраним в их связях несколько образов --- наборов по $pN$ одновременно активных нейронов --- правилом Хебба~\cite{Hebb1949}. Если подать на такую сеть половину образа, связи возбудят недостающую половину: сеть \emph{дописывает} образ по части, как группа на рис.~\ref{fig:bistable} поддерживала сама себя. Это ассоциативная память: адресом служит часть содержимого. \nt{GABA}, как в главе~\ref{sec:gaba}, держит общее число активных нейронов около $pN$, чтобы вместе не могли зажечься два образа.

Обозначим уровень ацетилхолина $a$, от $0$ до $1$, и запишем три его действия в модель прямо:
\begin{empheq}[box=\keybox]{equation}
\tau\frac{\dd r_i}{\dd t} = -r_i + F\Bigl(\tfrac{1+a}{2}\,x_i + (1-a)\sum_j W_{ij}\,r_j - g\Bigr),
\qquad
\frac{\dd W_{ij}}{\dd t} = \eta\,a\,(r_i-p)(r_j-p) .
\label{eq:acetnet}
\end{empheq}
Здесь $r_i$ --- частота нейрона $i$, $x_i$ --- его вход от органов чувств, $F$ --- пороговая передаточная характеристика, $g$ --- общее торможение от \nt{GABA}, растущее, когда активных нейронов больше $pN$. Множитель $\frac{1+a}{2}$ --- усиление входа, $1-a$ --- ослабление внутренних связей, $\eta a$ --- скорость обучения. Второе уравнение --- правило Хебба~\eqref{eq:hebb} в форме, удобной для редких образов~\cite{Tsodyks1988}: вес растёт, когда оба нейрона активнее обычного, и уменьшается, когда активен только один. Отрицательная часть весов, как всегда, стоит за \nt{GABA}-посредниками. Часов нет: и нейроны, и веса меняются непрерывно.

\begin{figure}[t]
\centering
\begin{tikzpicture}[>=Stealth,
  nrn/.style={circle,draw,very thick,fill=blue!6,minimum size=0.8cm,inner sep=0pt,font=\small},
  acet/.style={circle,draw,very thick,double,double distance=1.2pt,fill=orange!15,minimum size=1.35cm,inner sep=0pt,font=\scriptsize},
  bc/.style={->,thick,densely dashed,orange!85!black}]
\foreach \i/\y in {1/2.4,2/1.2,3/0}{
  \node[nrn] (N\i) at (3,\y) {$r_\i$};
  \draw[->,thick,blue!60!black] (0,\y) node[left,black] {\small $x_\i$} -- (N\i);}
\node[left,align=right,font=\small] at (-0.5,1.2) {от органов\\чувств};
\draw[->,thick,gray!60!black] (N1) to[bend left=30] (N2);
\draw[->,thick,gray!60!black] (N2) to[bend left=30] (N1);
\draw[->,thick,gray!60!black] (N2) to[bend left=30] (N3);
\draw[->,thick,gray!60!black] (N3) to[bend left=30] (N2);
\draw[->,thick,gray!60!black] (N1.east) .. controls (4.6,2.0) and (4.6,0.4) .. (N3.east);
\node[right,font=\small,gray!40!black,align=left] at (4.7,1.2) {внутренние\\связи $W$};
\node[acet] (A) at (1.2,-1.9) {\nt{ACET}};
\draw[bc] (A) -- (1.6,-0.15);
\node[font=\small,orange!60!black,anchor=east] at (0.95,-0.9) {$+$ вход};
\draw[bc] (A) -- (4.3,0.6);
\node[font=\small,orange!60!black,anchor=west] at (4.3,0.1) {$-$ внутренние связи, $+$ скорость обучения};
\node[right,font=\small,align=left] at (2.2,-1.9) {высокий уровень: «писать»\\низкий уровень: «читать»};
\end{tikzpicture}
\caption{\nt{ACET} как линия разрешения записи. Нейроны $r_i$ получают вход $x_i$ от органов чувств (синие стрелки) и возбуждают друг друга через внутренние связи $W$ (серые), в которых хранятся образы. Ацетилхолин (штриховые стрелки) усиливает вход, ослабляет внутренние связи и облегчает изменение весов, как в~\eqref{eq:acetnet}.}
\label{fig:acetnet}
\end{figure}

\section{Запись и чтение мешают друг другу}

Проверим модель~\eqref{eq:acetnet} на задаче, в которой запись и чтение действительно конфликтуют. Сеть из $200$ нейронов уже хранит пять образов по $20$ активных нейронов. Теперь нужно запомнить шестой, который наполовину совпадает с одним из старых: десять нейронов у них общие. Так бывает постоянно: новая комната похожа на знакомую, новое лицо --- на старое.

\emph{Запись при низком $a$.} Сначала отделим первые два действия ацетилхолина от третьего: оставим скорость обучения полной, а меняться будут только усиление входа и ослабление внутренних связей. При $a=0$ вход зажигает общие десять нейронов, а они через внутренние связи зажигают оставшуюся половину \emph{старого} образа. \nt{GABA} не даёт гореть обоим сразу, и старый образ, поддержанный собственными связями, вытесняет новый, который держится только на входе. Сеть видит не то, что перед ней, а то, что помнит, --- и правило Хебба старательно записывает увиденное.

В итоге новый образ вообще не запомнился: по его отличительной половине сеть ничего не вспоминает. А старый испорчен: вызванный своей отличительной половиной, он тянет за собой в среднем шесть чужих нейронов из десяти. При $a=0.1$ новый образ уже запоминается, но сливается со старым, и порча даже сильнее: каждый из двух, вызванный своей половиной, тянет за собой около восьми чужих нейронов; начиная с $a=0.2$ запись чистая: лишних нейронов при вспоминании меньше одного (усреднение по десяти прогонам).

\emph{Запись с настоящей скоростью обучения.} Теперь пусть ацетилхолин, как в~\eqref{eq:acetnet}, задаёт и скорость обучения. За одно предъявление новый образ запоминается целиком только при $a\ge0.9$; при $a=0.75$ --- на восемь десятых, при $a\le0.5$ не запоминается совсем.

\emph{Чтение.} Подадим на сеть с пятью чисто записанными образами половину образа, а потом уберём и её. При $a\le0.2$ сеть дописывает образ целиком и держит его сама; при $a=0.3$ --- почти целиком; при $a\ge0.4$ внутренние связи слишком ослаблены, и после того как подсказка исчезла, активность гаснет (рис.~\ref{fig:acetsim}).

\begin{figure}[t]
\centering
\begin{tikzpicture}[>=Stealth,x=9cm,y=3.2cm]
\draw[->,gray!70!black] (0,0) -- (1.1,0) node[right,black] {\small $a$};
\draw[->,gray!70!black] (0,0) -- (0,1.2);
\foreach \x in {0,0.2,0.4,0.6,0.8,1.0}{\draw[gray!60!black] (\x,-0.02) -- (\x,0.02); \node[below,font=\scriptsize] at (\x,0) {$\x$};}
\foreach \y in {0,0.5,1}{\draw[gray!60!black] (-0.01,\y) -- (0.01,\y); \node[left,font=\scriptsize] at (0,\y) {$\y$};}
\node[rotate=90,font=\small] at (-0.1,0.6) {доля восстановленного образа};
\draw[very thick,blue!60!black] plot[mark=*,mark size=1.6pt] coordinates {(0,1.00) (0.1,1.00) (0.2,1.00) (0.3,0.96) (0.4,0.04) (0.5,0.00) (0.75,0.00) (1.0,0.00)};
\draw[very thick,red!70!black] plot[mark=*,mark size=1.6pt] coordinates {(0.1,0.00) (0.2,0.00) (0.3,0.00) (0.5,0.00) (0.75,0.80) (0.9,1.00) (1.0,1.00)};
\node[font=\small,blue!60!black,anchor=west,align=left] at (0.03,0.5) {чтение:\\образ\\по половине};
\node[font=\small,red!70!black,anchor=east,align=right] at (1.0,0.35) {запись:\\новый образ\\за раз};
\end{tikzpicture}
\caption{Модель~\eqref{eq:acetnet}: $200$ нейронов, образы по $20$ активных, среднее по десяти прогонам. Синие точки --- чтение: какая доля образа восстановлена по его половине после того, как подсказка убрана. Красные точки --- запись: какая доля нового образа, наполовину похожего на старый, вспоминается после одного предъявления, если ацетилхолин задаёт и скорость обучения. Чтение работает при $a\le0.3$, запись --- при $a\ge0.9$.}
\label{fig:acetsim}
\end{figure}

\begin{keyblock}
\begin{proposition}[Запись и чтение требуют разных режимов]
\label{prop:writeread}
В модели~\eqref{eq:acetnet}, где одни и те же связи хранят старое и учатся новому, а ацетилхолин задаёт и скорость обучения, нет уровня $a$, при котором одинаково хорошо идут и запись, и чтение. Для записи внутренние связи нужно ослабить, иначе сеть запишет не вход, а вспомненное; для чтения их нужно усилить, иначе она не допишет образ по части. Нужен переключатель, и один широковещательный провод может им служить.
\end{proposition}
\end{keyblock}

Если бы ацетилхолин не менял скорость обучения, для обоих дел годилась бы узкая полоса $0.2\le a\le0.3$. Но тогда сеть училась бы и при каждом чтении, то есть заново записывала бы прочитанное вместе с его ошибками. Сама идея --- подавлять внутренние связи во время обучения --- взята из моделей, построенных по измерениям на срезах~\cite{HasselmoAndersonBower1992}. Числа выше относятся к нашей упрощённой модели; где именно лежат границы режимов в мозге, неизвестно.

\section{Сколько верить глазам}

Переключатель «писать/читать» --- крайний случай более общего вопроса: насколько верить входу и насколько --- памяти? Для этого вопроса у инженеров есть точный ответ, и он проливает свет на то, почему один провод управляет сразу и режимом, и скоростью обучения.

Пусть сеть следит за величиной $s(t)$, которая медленно блуждает: за секунду её дисперсия растёт на $q$. Органы чувств сообщают $y(t)=s(t)+$шум с интенсивностью $R$. Лучшая оценка $\hat s$ в непрерывном времени даётся фильтром Калмана--Бьюси~\cite{KalmanBucy1961}:
\begin{empheq}[box=\keybox]{equation}
\frac{\dd\hat s}{\dd t} \;=\; \frac{P}{R}\,\bigl(y-\hat s\bigr),
\qquad
\frac{\dd P}{\dd t} \;=\; q-\frac{P^2}{R} ,
\label{eq:kalman}
\end{empheq}
где $P$ --- дисперсия ошибки оценки, то есть насколько сеть сама не уверена в том, что знает. Первое уравнение --- та же RC-цепь, что у мембраны в модели нейрона~\eqref{eq:glutneuron}, но с постоянной времени $R/P$, которая меняется. Когда сеть не уверена, $P$ велико, постоянная времени мала, и оценка быстро следует за входом: это «писать». Когда уверена, $P$ мало, и оценка почти не слушает вход, держась за память: это «читать».

В установившемся режиме $P=\sqrt{qR}$ и постоянная времени памяти равна $\sqrt{R/q}$: если мир уходит на единицу за сто секунд, $q=0.01$, а шум органов чувств $R=1$, память должна помнить около десяти секунд.

Главное здесь то, что одно число $P/R$ задаёт сразу две вещи: и вес входа по сравнению с памятью, и скорость, с которой меняется хранимая оценка. Это ровно то, что делает ацетилхолин в~\eqref{eq:acetnet}. По гипотезе~\cite{YuDayan2005}, ацетилхолин и сообщает сети величину, подобную $P$, --- \emph{ожидаемую} неопределённость, ту, о которой сеть знает заранее. Медленным этот канал бывает не всегда: часть \nt{ACET}-нейронов отвечает на награду и наказание за два десятка миллисекунд, тем сильнее, чем неожиданнее подкрепление~\cite{Hangya2015}.

\begin{keyblock}
\begin{proposition}[Одна ручка для веса входа и скорости обучения]
\label{prop:kalman}
В оптимальном фильтре~\eqref{eq:kalman} вес, с которым вход смешивается с памятью, и скорость обучения --- одна и та же величина $P/R$. Поэтому разумно управлять ими одним сигналом. При неизменных свойствах мира этот сигнал выходит на постоянный уровень $\sqrt{q/R}$, и подстраивать его нужно, только когда свойства мира меняются.
\end{proposition}
\end{keyblock}

Второе предложение утверждения объясняет результат главы~\ref{sec:nora}, где подстройка скорости обучения по неожиданности не выиграла у лучшей постоянной: в мире с неизменными свойствами перемен лучшая скорость обучения и есть постоянная. Выигрыш появляется, когда мир внезапно становится другим. Тогда оценка вдруг оказывается далеко от входа, а $P$ об этом не знает. Правильное действие --- резко поднять $P$ и тем самым перейти в режим записи. По гипотезе, которую мы обсуждали в главе~\ref{sec:nora}, о такой \emph{неожиданной} перемене сообщает \nt{NORA}~\cite{YuDayan2005}. Разделение труда получается естественным: \nt{NORA} замечает, что мир изменился, \nt{ACET} держит сеть в режиме записи, пока новая обстановка не выучена.

\section{Сон как режим чтения}

Если высокий уровень ацетилхолина --- режим записи, то что делает сеть при низком? В медленном сне уровень ацетилхолина падает, внутренние связи освобождаются от подавления, входа от органов чувств почти нет. Сеть в режиме чтения без подсказки проигрывает то, что в ней записано. Регистрация активности показывает, что нейроны, работавшие днём вместе, в медленном сне снова срабатывают вместе~\cite{WilsonMcNaughton1994}, и даже в том же порядке~\cite{Skaggs1996}. По гипотезе~\cite{Hasselmo1999}, эти повторы переписывают быстро выученное днём в долговременную память (искусственное удлинение коротких вспышек повторов улучшает память~\cite{FernandezRuiz2019}): днём --- запись с входа, ночью --- перезапись изнутри. Для инженера это похоже на запись в быстрый буфер днём и его сброс в медленную постоянную память ночью.

\section{Итог}

\begin{table}[t]
\centering
\footnotesize
\setlength{\tabcolsep}{4pt}
\renewcommand{\arraystretch}{1.15}
\begin{tabular}{>{\raggedright}p{3.2cm}>{\raggedright}p{4.2cm}>{\raggedright\arraybackslash}p{6.0cm}}
\toprule
Что делает \nt{ACET} & Как & Что известно \\
\midrule
ослабляет внутренние связи & множитель $1-a$ в~\eqref{eq:acetnet} & измерено на срезах \\
усиливает вход & множитель $\frac{1+a}{2}$ & повышение возбудимости и передачи по входным путям измерено \\
облегчает обучение & скорость $\eta a$ & измерено \\
переключает «писать/читать» & утверждение~\ref{prop:writeread} & следует из моделей; прямого измерения режимов нет \\
сообщает ожидаемую неопределённость & $P$ в~\eqref{eq:kalman} & гипотеза \\
освобождает сеть для повторов во сне & низкий $a$ в медленном сне & низкий уровень во сне и повторы измерены; перезапись в долговременную память --- гипотеза \\
\bottomrule
\end{tabular}
\caption{Что добавляет \nt{ACET} к сети из \nt{GLUT}, \nt{GABA}, \nt{DOPA} и \nt{NORA}.}
\label{tab:acet}
\end{table}

\nt{DOPA} говорит сети, куда менять веса, \nt{NORA} --- насколько решительно пользоваться тем, что она знает, а \nt{ACET} --- слушать ли ей мир или себя (таблица~\ref{tab:acet}). Это последняя широковещательная линия управления из таблицы~\ref{tab:types}: разрешение записи, без которого память, хранящая образы в тех же связях, которыми она их читает, портила бы себя при каждом чтении.

\section{Задачи}

\begin{exercise}[\lvlA]
\label{ex:09:1}
\emph{Сколько помнить.} Фильтр~\eqref{eq:kalman} с $q=0.01$, $R=1$ помнит около $10$~с. Найдите $P_\infty$ и память $\sqrt{R/q}$, если (а)~шум датчика вырос вдвое по амплитуде; (б)~мир стал блуждать в сто раз быстрее. (в)~Во сколько раз должен вырасти шум, чтобы сеть помнила минуту?
\end{exercise}

\begin{exercise}[\lvlB]
\label{ex:09:2}
\emph{Режим записи после перемены.} Мир изменился, и неопределённость фильтра~\eqref{eq:kalman} с $q=0.01$, $R=1$ подскочила до $P_0\to\infty$. (а)~С какой постоянной времени $P$ возвращается к $P_\infty$? (б)~Через сколько секунд скорость обучения превысит установившуюся не больше чем на $10\,\%$?
\end{exercise}

\begin{exercise}[\lvlB]
\label{ex:09:3}
\emph{Постоянная скорость обучения.} Сеть учится как $\dd\hat s/\dd t=(y-\hat s)/T$; мир блуждает, $\dd s/\dd t=w$, $y=s+n$, где $w$ и $n$ --- белые шумы с интенсивностями $q$ и $R$. Найдите лучшее $T$ и дисперсию ошибки при нём. Во сколько раз она вырастет, если $T$ ошибочно в $2$ и в $10$ раз?
\end{exercise}

\begin{exercise}[\lvlB]
\label{ex:09:4}
\emph{Где граница записи.} В \texttt{models/\allowbreak acet\_memory.py} порог $\theta=0.35$, вес внутри образа $w=0.041$ (в идеале $0.045$). Новый образ делит $m=10$ нейронов со старым. При каком $a$ в~\eqref{eq:acetnet} остальные нейроны старого образа не зажигаются от этих $m$ (порог резкий)?
\end{exercise}

\begin{exercise}[\lvlB]
\label{ex:09:5}
\emph{Моделирование: ёмкость памяти.} Измените в \texttt{models/\allowbreak acet\_memory.py} функцию \texttt{read\_sweep}: при $a=1$ запишите $M$ образов ($M$ от $5$ до $100$), затем при $a=0$ прочитайте первые десять по половине. При каком $M$ появляются ошибки и как предельное $M_{\max}$ зависит от $N$ и $p$?
\end{exercise}

\begin{exercise}[\lvlB]
\label{ex:09:6}
\emph{Проектирование: запись без утечки.} Со скоростью обучения $\eta a$ сеть учится и при чтении. Предложите монотонную $\eta(a)\le\eta$, равную нулю при $a\le0.3$ и не хуже $\eta a$ при $a\ge0.9$. Проверьте: после $20$ чтений при $a=0.2$ с подсказкой, где три чужих нейрона, сколько их вошло в образ?
\end{exercise}

\begin{exercise}[\lvlC]
\label{ex:09:7}
\emph{Как переписать буфер ночью.} (а)~Во сне $a=0.05$, повтор длится $0.1\,\text{с}$ против $0.2\,\text{с}$ днём при $a=1$. Сколько повторов накопят изменение весов одного дневного предъявления при скорости $\eta a$ и при $\eta(a)$ задачи~\ref{ex:09:6}? (б)~Хранилище учится в $100$ раз медленнее буфера и слышит образ $20$ раз за ночь по $0.1\,\text{с}$. Сколько нужно ночей?
\end{exercise}
